\section{Source-authenticated crushing and positive quadrilateral search}
\label{sec:normal-crushing-search}

The preceding filtered search stayed on one triangulation and accumulated
rejected-support clauses.  The new unmarked route consumes a positive normal
ray by crushing and restarts on a strictly smaller triangulation.  This
removes the need to bound the number of geometric rejections on a fixed
source.  It also changes the required discovery query: any positive-Euler
non-link ray suffices, so quadrilateral corners suffice even when they do
not include an essential disc.  The implementation provides independently
replayed quotients, diagram-bound successor chains, and a complete additive
recognition backend.  General quasi-polynomial discovery remains open.

\subsection{The topological premise and its source authentication}

We use the classical crushing theorem in the form of
Burton--Ozlen~\cite[Observations~3--4 and Lemma~5]{burton-ozlen}.
For a classical knot exterior, crushing a connected positive-Euler normal
surface with quadrilateral content either leaves the same exterior as the
unique torus-boundary component or proves that the knot is trivial when no
such component remains.  The surviving tetrahedra are precisely those with
zero quadrilateral coordinates.  Thus a retained successor is strictly
smaller.  The no-torus conclusion also holds at one tetrahedron: the
$t>1$ hypothesis in Lemma~5 permits replacing that case by a smaller
one-tetrahedron solid torus, whereas our checker terminates directly.

The low-level source validator establishes a finite compact orientable
manifold with one torus boundary.  It does not establish that the manifold
is a knot exterior.  Accordingly, \path{normal_crushing.py} produces only
a conditional quotient packet.  In \path{normal_crushing_chain.py}, the
independent canonical diagram-exterior checker authenticates the initial
source; optional boundary shellings are independently replayed; every
subsequent quotient is authenticated against the preceding selected source.
Only this diagram-bound chain interprets absence of a torus as an unknot
certificate.  A partial chain with status \texttt{REDUCED} authenticates a
smaller exterior and has no recognition verdict.

This distinction is necessary.  A compact orientable torus-boundary
manifold may have a closed connected-sum factor that crushing removes.
Source-format validation alone would then be insufficient to identify a
surviving component with the original manifold.  The chain's inductive
knot-exterior premise supplies the missing topological information.

\subsection{Primitive full-source rays supply connectedness}

Coordinates are validated against every original normal matching equation,
nonnegativity and per-tetrahedron quadrilateral compatibility.  The producer
restricts the original matching matrix to the positive support and checks
nullity one.  The verifier reconstructs those columns and uses its own
reverse-pivot kernel chart to check primitive integral scale and full-source
extremality.  This one-shot check avoids preparing all $7t$ columns merely
to validate a single ray.  Euler characteristic is computed from exact face,
edge and normal-disc counts, without expanding any coordinate multiplicity.
A positive Euler value and at least one quadrilateral are mandatory.
Explicitly, for $S=\operatorname{supp}(x)$ the rank condition is
$\operatorname{rank}(A_{[:,S]})=|S|-1$, using exact rational elimination.

\begin{lemma}[Connected primitive normal ray]
An admissible primitive integral extreme ray of the full matching cone
represents a connected normal surface.
\end{lemma}
\begin{proof}
Every connected component has a nonzero nonnegative integral matching
vector, and these vectors sum to the supplied vector $x$.  Extremality
makes each component vector $\lambda_i x$ for some $\lambda_i>0$.
Since the coordinates of $x$ have greatest common divisor one, an integer
linear combination of them equals one.  Applying that combination to
$\lambda_i x$ shows that $\lambda_i$ is an integer.  The equality
$\sum_i\lambda_i=1$ therefore permits only one component.
\end{proof}

The full matching cone matters.  An arbitrary extra linear constraint can
make a disconnected vector extreme in a smaller cone.  Coordinate-zero
sector restrictions are faces and do not cause this problem, but the
crushing checker always retains every original matching equation and adds
no new linear constraints.
In an authenticated knot exterior, a connected positive-Euler embedded
surface is a sphere or disc, since an embedded projective plane is excluded.
Any quadrilateral excludes a vertex link and guarantees strict size descent.
Triangle-only, nonprimitive, nonpositive and nonextreme inputs are rejected
by this ray interface rather than silently normalized or repurposed.

\subsection{Face states, erased cycles and independent quotient replay}

For a tetrahedron with nonzero quadrilateral type $q$, the two vertex pairs
are $\{0,q+1\}$ and its complement.  The local face transition exchanges
the two vertices in the pair containing the entry face label and fixes
the remaining pair.  The producer composes these local permutations with
source face gluings, starting from every retained face.  A path ends at
another retained face or at the source boundary.  The resulting endpoint
and permutation specify the quotient face record.

The directed state is a tetrahedron together with its entry face.  A path
may revisit a tetrahedron through a different face, so a tetrahedron-only
cycle guard is incorrect.  Distinct retained-face traversals use disjoint
directed entry states; reverse traversal uses the paired exit states.
There are at most $4t$ such states.  This bounds the producer's aggregate
face-transition work independently of the normal coordinate magnitudes.

The independent checker does not call this tracer.  It builds a labelled
graph on the $4t$ face states.  Source face pairings and internal quad-pair
transitions are its edges.  Each vertex has degree at most two.  A connected
component with two retained endpoints determines a quotient gluing, obtained
by propagating permutation labels.  A component with one retained endpoint
ends at the boundary.  Components with no retained endpoint are erased and
may be cycles.  Such isolated cycles are legitimate; only a repeated state
reachable from a retained endpoint contradicts the path structure.

The packet contains original coordinates, retained source indices, the
entire quotient, its complete component census, the chosen torus component,
and the renumbered successor.  The checker derives each field independently.
Every quotient component is validated, including components discarded from
the successor: face pairings must be reciprocal, edge reversals are excluded,
orientations must be consistent, and vertex links must be finite spheres or
discs.  Boundary components are assembled explicitly and their Euler
characteristics computed.  Selecting the unique component with one torus
boundary therefore does not confuse a closed component or a failed manifold
validation with disappearance of the torus.

The component census is shared validated topology code; quotient discovery
and source extremality have separate producer and checker implementations.
Regina is used only as an external research oracle, never as a dependency
of a maintained producer or certificate verifier.  Callbacks are guarded
through nested verification so that caller exceptions cannot be swallowed
as malformed data or converted into resource exhaustion.

\subsection{Why positive quadrilateral corners are now enough}

Fix one allowed quadrilateral type per tetrahedron, permitting its coordinate
to be zero.  Let $Q$ be the nonnegative quadrilateral matching cone.  The
triangle matching equations supply linear potentials $\ell_{vc}(q)$ at
the corners incident to each global vertex $v$.  The canonical lift subtracts
their minimum at that vertex, making at least one triangle coordinate zero.
As established in the all-dimensional Euler-convexity result, its Euler
characteristic has the form
\[
 F(q)=L(q)-\sum_v\chi(V_v)\min_c\ell_{vc}(q),
\]
where $L$ is linear and each vertex link $V_v$ has Euler value one or two.
Consequently $F$ is convex and positively homogeneous.

\begin{proposition}[A positive Q corner suffices for crushing]
If this sector contains a canonical surface with positive Euler
characteristic and nonzero quadrilateral vector, it contains a positive
canonical lift of a Q-cone extreme ray.  The primitive lift is a full
matching-cone extreme ray and is eligible for the crushing interface.
\end{proposition}
\begin{proof}
Intersect $Q$ with $\sum_jq_j=1$.  This is a compact polytope when nonempty.
Writing a point as a convex combination of its vertices, convexity bounds
its $F$ value above by the corresponding weighted average of vertex values.
Thus a positive value forces a positive vertex value.  Homogeneity preserves
the sign when a ray is made primitive.

Let $x$ be a canonical lift of such an extreme Q ray and suppose $x=u+w$
in the nonnegative full matching cone.  Projecting to Q gives proportional
vectors, say $q(u)=\lambda q(x)$ and $q(w)=(1-\lambda)q(x)$, with
$0\leq\lambda\leq1$.  The difference $u-\lambda x$ is triangle-only
matching data, hence a linear combination of vertex links.  At each vertex,
a zero triangle coordinate of $x$ is also zero in $u$ and $w$, so the
corresponding link coefficient in the difference is zero.  Therefore
$u=\lambda x$ and $w=(1-\lambda)x$.  This proves full-cone extremality.
The preceding lemma proves connectedness of its primitive representative.
\end{proof}

For direct essential-disc selection a positive Q corner may be inadequate:
it can have trivial boundary while an essential standard ray lies elsewhere.
Crushing makes that distinction unnecessary for the next progress step.
The new driver therefore uses only the quadrilateral phase of the native
sector enumerator and restarts after its first positive ray.  It does not
build standard minimum-envelope arrangements or support-exclusion masks.

\subsection{Complete global search and independently checked negative coverage}

At a source with $t$ tetrahedra, enumerate the $3^t$ maximal type assignments.
These sectors cover every admissible normal vector: tetrahedra with no
quadrilateral can be assigned any type.  In each sector enumerate all Q
corners until a positive lift is found.  Crush and independently replay that
lift, append the source-bound step, and restart on the selected successor.
No torus ends with an independently checked \texttt{UNKNOT} certificate.
If every sector completes without a positive lift, return \texttt{KNOTTED}
with the complete final-source sector exhaustion packets.

The producer retains the existing sufficient matching-row screen for
source-proven forced-zero types when the initial Q nullity exceeds three.
It searches the resulting face and zero-expands negative witnesses back
into the full original assignment.  The independent exhaustion checker
replays any accompanying matching-support proof against that original source.
This screen is polynomial but need not identify every forced zero.

The default traversal first draws sixteen assignments with the fixed seed
37, discarding duplicates, and then visits all remaining assignments in
lexicographic order.  This bounded deterministic prefix is a heuristic:
neither its effectiveness nor completeness relies on a probability model.
The \texttt{canonical} option omits the prefix.  Both orders visit each
assignment exactly once per source, and completed negative packets are
sorted into canonical order before their independent replay.

\begin{proposition}[Completeness of unbounded crushing search]
On a classical knot diagram, unbounded search terminates with the correct
verdict.  There are at most $t_0$ crushing steps, and at most $t_0-1$ steps
retaining a nonempty successor, where $t_0$ is the initial tetrahedron count.
\end{proposition}
\begin{proof}
The diagram and shelling checkers establish the initial exterior.  Each
verified nonterminal crush preserves that invariant and strictly decreases
the number of tetrahedra; a terminal crush proves triviality by the stated
topological theorem.  There cannot be more than $t_0$ strict decreases to
zero.  Every individual sector search is finite.

If the current exterior belongs to the unknot, the normal-disc theorem gives
an essential normal disc.  Subtractable vertex links can be realized as
disjoint components, so this connected disc has zero minimum triangle
multiplicity at every vertex.  It has no vertex-link component, so its normal
vector is canonical, and it has nonzero quadrilateral content.  Its sector
therefore contains a positive Q corner by the preceding proposition.  Hence
complete exhaustion of every sector cannot occur for the unknot.  Conversely,
a verified no-torus terminal cannot occur for a nontrivial knot.  Finiteness
of all searches and strict descent leave exactly the asserted outcomes.
\end{proof}

The recognition certificate verifier imports no discovery routines.  For
an unknot it checks the whole chain and requires no negative-sector packet.
For a knottedness claim it first authenticates the final reduced exterior,
then requires every maximal assignment exactly once in canonical order,
with a completed nonpositive Q-sector certificate.  The existing independent
dense sector checker reconstructs the Q corners and their Euler values.
A missing sector, a partial search, or an unsuccessful checker call cannot
establish the negative verdict.

Configured sector and candidate limits are shared across all sources and
stages; a separate crushing limit bounds performed successor steps.  Reaching
a limit returns \texttt{INCONCLUSIVE} with an authenticated partial chain.
Preparation is not charged as a candidate, and these counts are not total
time or memory budgets.  Caller cancellation remains distinct from the
backend's own limits.  This complete backend is additive; the default
recognizer's dispatch policy is unchanged.

\subsection{Asymptotic accounting and the remaining target}

In a sector with Q nullity $d\geq1$ and at most $t$ allowed types, both native
Q enumeration and independent dense Q replay solve at most
$\binom{t}{d-1}\leq2^t$ candidate systems.  Nullity zero has no Q rays
and requires no candidates.  Each system has polynomial exact-bit
cost.  Canonical primitive coordinates satisfy the inherited exponential
height bound and thus have $O(t)$ bits.  Combining $3^t$ assignments and at
most $t_0$ decreasing stages gives a crude $6^{t_0}\operatorname{poly}(t_0)$
algebraic-work bound.  Current Python hash tables do not establish that
literal comparison-model bound against adversarial collisions; even charging
quadratic table work preserves the unconditional $2^{O(t_0)}$ bit-time class.
Neither bound is quasi-polynomial in diagram size.

The useful conditional statement is sharper than a bound on fixed-source
rejection counts.  Suppose a complete positive-nonlink oracle can be built
and run in quasi-polynomial time on every successor of binary size polynomial
in the initial diagram size.  The polynomial number of crushing stages and
their polynomial source/surface/quotient checks preserve quasi-polynomial
time.  No marked annulus, suture or peripheral-cochain transport is required
for this unmarked route.  The missing condition is now a complete efficient
positive search on every smaller exterior, not a geometric bound on the
support-filtered algorithm's rejected-ray count.  The present exhaustive
sector family does not satisfy that condition.
